% !PRINTED 341
\weberpage{341}
%<<CONT
ist, so lange $\xi, \eta$ dem Bereich (8) angehören. Dann folgt aus (7)
\[
\beta' - \alpha' < (\beta - \alpha)\Theta.
\]
Ebenso folgt, wenn wir auf dieselbe Weise von dem Intervall $(\alpha', \beta')$ zu einem engeren Intervall $(\alpha'', \beta'')$ fortschreiten,
\[
\beta'' - \alpha'' < (\beta' - \alpha')\Theta',
\]
worin $\Theta'$ dieselbe Bedeutung für $\alpha', \beta'$ hat, wie $\Theta$ für $\alpha, \beta$. Da aber $\alpha', \beta'$ dem Bereich (8) angehören, so ist $\Theta'$ nicht grösser als $\Theta$ und statt $\Theta'$ kann auch $\Theta$ gesetzt werden. Es folgt also:
\[
\beta'' - \alpha'' < (\beta - \alpha)\Theta^2,
\]
und so schliessen wir weiter
\[
\beta^{(\nu)} - \alpha^{(\nu)} < (\beta - \alpha)\Theta^\nu.
\]

Die Intervalle nehmen also mindestens so stark ab, wie die Glieder einer fallenden geometrischen Progression.

Als Beispiel mag die Gleichung dienen:
\[
x^3 - 2x^2 - 2 = 0,
\]
die eine Wurzel zwischen
\[
\alpha = 2{,}35 \quad \text{und} \quad \beta = 2{,}36
\]
hat.

Man erhält aus den Formeln (4) für
\[
\beta' = 2{,}35931\ldots, \quad \alpha' = 2{,}359298\ldots,
\]
so dass ein in der vierten Decimale genauer Werth der Wurzel
\[
2{,}3593
\]
ist. Für diesen Werth selbst ist, wie eine genauere Rechnung ergiebt, $f(x)$ noch negativ, so dass er für $\alpha'$ genommen werden kann. Der nächste Schritt der Annäherung ergiebt.
\[
2{,}359304.
\]

\begin{center}
\S.\,110.

\gesperrt{Die Näherungsmethode von Daniel Bernoulli und verwandte Methoden.}
\end{center}

Die Methode zur genäherten Auflösung einer Gleichung, die von \gesperrt{Daniel Bernoulli} herrührt, beruht darauf, dass, wenn man eine Reihe reeller Grössen hat, die Potenzen der grössten unter ihnen um so mehr die gleich hohen Potenzen der übrigen überwiegen werden, je höher die Potenzen sind.
